\section{选模型族就是选一组关于世界形状的假设}
\label{sec:18-Synthesis/01-System-Lifecycle/03}

一家连锁品牌要做门店周销量预测。手上的数据是两百家门店、十八个月，合起来三千六百行，每行四十来个特征：门店面积、周边人口、天气、促销标记、节假日、上周销量。

方案评审会上出现了两派。一派要上序列模型，理由是销量本来就是时间序列，注意力能自动学到长程依赖；另一派要用梯度提升树，理由是``这种表格数据树就够了''。争论的形式是``哪个模型更强''，而这个问题问错了。

后来两个方案都做了。序列模型的验证误差比树高出约三成，怎么调都追不上。原因不复杂：三千六百行数据里真正独立的信息量远小于这个数（同一家门店相邻两周高度相关），而注意力几乎不对函数形状做任何限制，它要靠数据把``销量与面积大致单调''这种事实一点点学出来；树只需要几次分裂就把它表达了。

反过来的例子同样常见。同一家公司的图像质检项目上，有人先用梯度提升树在像素上训了一版，效果差得离谱——树的每次分裂只看一个像素，而``这张图里有没有划痕''这件事跟任何单个像素的取值都没什么关系。

这两次失败方向相反，根子是同一件事：\textbf{每一个模型族都在替你断言数据应当长成什么形状，而这个断言写错了不会报错，只会让结论悄悄有偏。}

\subsection{每个模型族断言了什么}

线性模型断言效应可加：各个因素各自贡献一份，合起来就是结果。\hyperref[sec:13-Machine-Learning/02-Linear-Models/01]{``线性回归从投影到概率模型''}给出的两种读法都指向这一点——几何上它把响应投影到特征张成的子空间，概率上它假设噪声独立同分布且与特征无关。假设成立时，系数直接可读、不确定性可量化、失效方式已知；假设不成立时（比如促销效果依赖于门店大小），模型不会抗议，它只会把交互项的影响摊派到各个主效应上去。

树断言决策边界是轴对齐的分段常数。\hyperref[sec:13-Machine-Learning/03-Trees-Ensembles-and-Neighbors/01]{``决策树把预测写成递归分区''}讲的正是这个结构：预测值只取决于样本落在哪个矩形块里。这个假设对表格数据出奇地合适，因为业务规则本来就常常是阈值形式的；对一条 \(45\) 度的斜边界则要用无数级台阶去逼近，而每一级都要花掉样本。

最近邻把假设完全交给度量。\hyperref[sec:13-Machine-Learning/03-Trees-Ensembles-and-Neighbors/03]{``KNN 把度量直接变成归纳偏置''}说得直白：模型本身几乎没有参数，全部先验都写在``什么算相近''里。这与上一篇的结论接上了——特征怎样缩放，就是在决定这个模型相信什么。

核方法把假设写进核函数。\hyperref[sec:13-Machine-Learning/09-Kernel-Methods-and-Gaussian-Processes/05]{``Gaussian process 把核变成函数先验''}把这件事说到底：选一个核等于选一个关于函数光滑程度、周期性、长度尺度的先验分布。用 RBF 核就是断言``输入接近则输出接近''，而``接近''的尺度由那个带宽参数定死。

概率图把假设写成``哪些依赖不存在''。\hyperref[sec:13-Machine-Learning/07-Graphical-and-Sequential-Models/01]{``图结构把联合分布拆成局部因子''}里，缺一条边比多一条边说得更多——它断言在给定某组变量之后，另外两个变量再无关联。这类假设最省参数，也最容易被现实推翻。

卷积断言平移等变与局部性：一个模式出现在图像的哪个位置不改变它的含义，而判断一个位置需要的信息集中在它附近。\hyperref[sec:14-Deep-Learning/04-Convolutional-Networks/01]{``卷积把平移结构写进线性映射''}说明这两条假设怎样把一个全连接层的参数量压到极小。注意力则几乎放弃了位置假设，允许任意两个位置直接对话，于是它连顺序都不知道——\hyperref[sec:14-Deep-Learning/05-Attention-and-Sequence-Models/03]{``位置必须外挂，因为注意力看不见顺序''}讲的就是这件事：把结构假设拿掉之后，得手工再加回去一部分。

\subsection{偏置强度与样本量是一条权衡曲线}

上面这些假设有强有弱，而强弱本身不分好坏，它只有在与样本量配对之后才有意义。

\begin{center}
\begin{tikzpicture}[x=1cm,y=1cm,font=\small,>=stealth]
  \draw[->,black!60] (0,0) -- (11.6,0) node[anchor=north east,font=\footnotesize,black!70] {样本量（对数刻度）};
  \draw[->,black!60] (0,0) -- (0,4.6);
  \node[rotate=90,anchor=south,font=\footnotesize,black!70] at (-0.45,2.3) {测试误差};
  \draw[dashed,black!40] (0,1.36) -- (11.2,1.36);
  \node[anchor=east,font=\footnotesize,black!60] at (11.2,1.72) {强偏置的天花板};
  \draw[very thick,blue!65!black,smooth] plot coordinates
    {(0.3,3.6) (1.5,2.4) (3,1.8) (4.5,1.52) (6,1.43) (8,1.38) (10.9,1.36)};
  \draw[very thick,orange!80!black,smooth] plot coordinates
    {(0.3,4.35) (1.5,4.0) (3,3.35) (4.5,2.6) (6,1.92) (7.4,1.36) (8.8,1.0) (10.9,0.6)};
  \draw[dashed,black!45] (7.4,0) -- (7.4,1.36);
  \node[font=\footnotesize,black!70,anchor=north] at (7.4,-0.08) {交叉点};
  \node[anchor=west,blue!65!black,font=\footnotesize] at (2.6,1.05) {强偏置（线性、树、卷积）};
  \node[anchor=west,orange!80!black,font=\footnotesize] at (5.3,3.6) {弱偏置（大容量网络）};
  \node[align=center,font=\footnotesize,black!70] at (3.4,4.2) {数据少：\\假设替你补足了信息};
  \node[align=center,font=\footnotesize,black!70] at (9.9,2.6) {数据多：\\假设成了上限};
\end{tikzpicture}

\emph{两条曲线必然相交。强偏置在左半边靠假设省下样本，右半边则被同一个假设卡住；争论``哪个模型更强''等于争论应当站在交叉点的哪一侧，而这取决于手上有多少数据。}
\end{center}

蓝线之所以会平掉，是因为它的误差里有一块与样本量无关：真实规律不在这个模型族能表达的范围内，再多数据也补不上。这块是近似误差，在\hyperref[sec:14-Deep-Learning/01-Representation-and-Approximation/02]{``通用逼近定理说清了什么又没说什么''}里被说清了另一半——逼近定理保证足够宽的网络可以把这块压到任意小，但它不告诉你需要多少神经元、也不保证优化找得到那组参数。所以橙线在左边高不是因为表达力不够，是因为可选的函数太多而数据不足以在其中挑出对的那个。

这条曲线也解释了门店销量那个案例：三千六百行落在交叉点左边很远的地方。而图像质检落在另一侧的另一个意义上——像素上的树不是偏置太强，是偏置的\emph{方向}错了，轴对齐分区跟图像的平移结构完全不搭。\hyperref[sec:14-Deep-Learning/04-Convolutional-Networks/05]{``归纳偏置何时帮忙何时挡路''}讲的就是这个区分：偏置的方向对不对，与偏置的强弱是两个独立的问题。方向对了，强偏置在小数据上是纯赚；方向错了，偏置越强败得越彻底。

至于``可选的函数太多''该怎么量化，\hyperref[sec:13-Machine-Learning/10-Learning-Theory/03]{``VC 维把无限假设类变有限''}给出了一种回答：假设类的有效大小不由参数个数决定，而由它能打散多少个点决定。这个量在实践中几乎算不出来，但它提供的直觉是可靠的——控制泛化的是假设类的容量，而正则化、参数共享、结构约束都是在压这个容量。

\subsection{深度学习换了一笔账，不是免了这笔账}

前面那些模型族都把``用什么表示''交给人来设计：特征工程写在模型之外，模型只负责在给定表示上拟合。\hyperref[ch:122]{``深度学习：让表示、梯度与数据结构在复合函数中相遇''}做的事是把表示本身变成可优化的对象，于是不再需要手工告诉模型``划痕看起来是什么样''。

这不是白得的。换来的东西是表示自动化，付出的东西是几乎所有保证都退化成经验规律：没有全局最优的承诺，没有系数的可解释性，没有关于需要多少数据的可用估计，连``加深一层会不会更好''都只能试。前面那些模型族之所以仍然值得优先考虑，不是因为它们更精确，是因为它们坏掉的时候你知道它们为什么坏。

在深度学习内部，模型族的选择同样是假设的选择，只是假设写在架构里而不是写在函数形式里。卷积、循环、注意力对应三种关于依赖结构的断言，选错了照样表现为``怎么调都上不去''。这一点在\hyperref[sec:14-Deep-Learning/05-Attention-and-Sequence-Models/03]{``位置必须外挂，因为注意力看不见顺序''}那样的细节里看得最清楚：架构拿掉了什么，就得从别处补回来什么。

\subsection{这一步做错时，故障长什么样}

第一种形态是模型怎么调都卡在某个水平。换超参、加数据、换优化器，验证误差在小数点后第三位上晃动。这几乎总是撞到了天花板——真实规律不在这个族里，而调参改变不了族。此时应当做的不是继续调，是换一个方向的假设，或者手工加入能表达缺失结构的特征（把交互项显式写出来往往立刻见效，而这本身就是在证实假设不对）。

第二种形态是训练集上完美，测试集上崩溃，而且加正则之后两边一起变差。这是偏置太弱且数据不足，模型有足够多的自由度记住训练数据。区分它与第一种的办法很简单：看训练误差。训练误差就降不下去是第一种，训练误差接近零而测试误差高是第二种。这两种病的药方相反，混淆它们会导致该换模型族的时候去加数据、该加数据的时候去调优化器。

第三种形态最容易被误读：模型在整体上还行，但在某一类输入上给出明显违反常识的输出（销量对面积非单调、图像稍微旋转一点结论就变）。这说明模型族的假设与真实结构在某个方向上不一致，而整体指标把这个不一致平均掉了。查它的办法是构造受控的输入变化，看输出是否按应有的方式响应。

模型族定下来之后，剩下的看起来只是把参数解出来。下一篇要说明的是，这一步提供的保证比大多数人以为的要少得多。
